
\begin{frame}{Contribution}
\hypertarget{src:contribution}{}
\small
Where a large literature reads courts as a \pos{leveling} force that opens contests to under-represented entrants, the evidence here documents the opposite margin: in single-winner municipal races, an expanded judicial arena \textbf{re-concentrates} the vote toward the front-runner without reshaping who runs or who wins --- and, where an incumbent is defended, pushes marginal voters off the ballot, the \negt{barrier} signature.
\medskip

The shift-share design isolates this from the nation-wide litigation trends common to every race, and the sharp \textbf{mayoral-versus-council asymmetry} shows the mechanism is specific to the single-winner contest rather than a general effect of legal activity.
\medskip

The reframing matters as much as the sign. The relevant margin of judicialization is its \emph{expanded role} --- rule-making, weaponization, and the anticipation of both --- of which observed lawsuit counts are only the visible tip. The effect operates even as raw filings hold flat or fall.
\end{frame}
